\subsection{What was already on disk, and what was added under predeclaration}

The study has two layers. The first is a re-reading of records written by runs
that finished before this analysis began: three comparisons, none of which was
generated for this paper. The second is two computations made for this paper
under written predeclarations, each frozen and hashed before its first result
existed: a perceptual scoring pass over the subset crops, and a
three-initialisation ablation of the toy transformer. Every quantity in both
layers is tied to the manuscript by an evidence-manifest binding; the figure
code verifies each bound file before reading it, re-derives the summaries from
the per-crop or per-cell rows, holds each result to its own predeclaration, and
stops the build on any mismatch.

Three comparisons are read. The first is a toy-scale training run in which a
small convolutional restoration network and a small transformer of comparable
size were trained under a matched step budget on the same data, evaluated on the
same \ToyN{} held-out pairs at the same number of sampling steps, and scored
against the same degraded inputs. The second is the same convolutional arm
retrained at \BudgetMultiple{} times the budget, written outside the results tree
and never promoted. The third is a released diffusion restoration model run over
the \SubsetN{} crops of the published RealSR crop128 protocol, scored against
reference images, alongside bicubic and nearest-neighbour upsampling of the same
inputs.

\subsection{The trivial arms, and why they are arms}

Each comparison has an arm that performs no restoration. At toy scale it is the
degraded input passed through unchanged, which the training record scores
alongside the two trained networks. At subset scale there are two: bicubic and
nearest-neighbour upsampling of the same low-resolution inputs, scored by the
same metric pass over the same crops.

These are treated here as arms rather than as context. An arm that does nothing
is the weakest hypothesis a restoration comparison can be asked to beat, and a
method that does not beat it has not been shown to restore anything, whatever
its ranking against other methods. The three comparisons below differ in what
the trivial arm reveals, not in whether it should have been reported.

\subsection{What the fidelity measures are, and what they are not}

Two fidelity measures appear: peak signal-to-noise ratio, computed over the full
colour image and separately over the luma channel with a four-pixel border
removed, and structural similarity~\cite{wang2004ssim} under the same two
conventions. Both compare a restored image with a reference pixel by pixel or
window by window, and both reward outputs that stay close to the reference.

Neither is a measure of perceptual quality, and the difference matters for
everything reported at subset scale. Restoration methods that synthesise
plausible detail are known to score worse on distortion measures than methods
that produce a blurrier output closer to the conditional
mean~\cite{blau2018perception,ledig2017srgan}, and the trade-off is a property
of the estimator rather than a defect of any particular
model~\cite{blau2018perception}. A comparison with a distortion measure alone
can establish which output is closer to the reference; it cannot establish
which output is better, and a comparison with a perceptual measure alone cannot
either. This paper reports both, and treats their disagreement as the result.

\subsection{The perceptual endpoint}

The perceptual endpoint is the learned perceptual image patch similarity
(LPIPS), version 0.1, with the AlexNet backbone and its released linear
calibration head~\cite{zhang2018lpips}: lower is closer. Each restored output
and its reference are the full \(512\times512\) colour crop, scaled to
\([-1,1]\), with no border removed. The primary contrast for crop
\(i\in\mathcal{M}\) is \(\Delta^{\mathrm{LPIPS}}_i =
\mathrm{LPIPS}_{\mathrm{bicubic},i}-\mathrm{LPIPS}_{\mathrm{model},i}\), so a
positive value is a crop on which the model is perceptually closer. The
primary test is the exact two-sided sign test on the \SubsetN{} paired
contrasts at \(\alpha=0.05\), ties counted and excluded; a percentile bootstrap
of the mean contrast over 1,000 resamples and the per-camera win counts are
descriptive. The protocol, including these choices and the three possible
readings of the outcome, was written and hashed before any perceptual value was
computed, and the figure code re-derives the wins, the test and the mean from
the per-crop rows.

Two disclosures attach to this endpoint. First, provenance: the bound fidelity
table was produced by an earlier inference run on another machine, and its
restored pixels were not the ones scored here. The perceptual pass drew a fresh
sample of the released model on the same \SubsetN{} low-resolution inputs, at
sampling seed \LpipsSeed{} with \LpipsSteps{} sampling steps and the released
weights, and scored that sample. A predeclared anchor rule governs whether the
fresh sample may be read beside the bound ranking: its mean RGB PSNR must lie
within half a decibel of the bound record's, and it must lose to bicubic on at
least ninety of the \SubsetN{} crops; had either failed, the perceptual numbers
were to be reported as not comparable and no rewriting of the thesis was to
occur. The anchor passed with a shift of \AnchorDeltaDb{} decibels and
\AnchorLosses{} losses, and the bicubic arm recomputed for the pass reproduces
the bound bicubic rows to floating-point precision. The perceptual comparison
is therefore between the fresh sample and bicubic; the PSNR comparison in the
same figure uses the fresh sample as well, and the bound PSNR table is unchanged.

Second, the predeclaration was amended once. The predeclared scoring script
crashed while serialising its summary, after the per-crop values had been
computed and before any formal endpoint file was written. The amendment
replaced the script with a version that casts the offending values to plain
types --- the endpoints, pairing, test, threshold and anchor rule were not
changed --- and recorded the secondary backbone as not computed (below). The
primary result had by then been observed on a provisional side run that
imported the predeclared script's own functions; the amendment says so, and the
formal run under the amended protocol reproduced the provisional run exactly:
the same \LpipsWins{} wins, the same sign-test probability to the last digit,
the same anchor shift. Both the original predeclaration and the amendment are
bound and hashed with the result.

The protocol named LPIPS with the VGG16 backbone as a secondary robustness
check, never to be promoted. It was not computed: the cached VGG16 checkpoint
on the scoring machine was truncated, and the download route the workspace
permits was unavailable at the time of the run. Its status is recorded as such
in the bound summary, the per-crop rows carry null values for it, and nothing
is imputed for it. A distinct learned quality score that would have required a
further download was likewise omitted rather than substituted.

\subsection{The initialisation ablation}

The toy transformer was retrained under three initialisations with everything
else fixed: the library's default construction, which is the arm behind the
bound toy record; a standard transformer scheme applied to the same graph
without any zeroing step; and that scheme plus exact zeroing of the adaptive
normalisation gates and the final projection, which is the published
recipe~\cite{peebles2023dit}. Each was trained at \InitSeeds{} seeds for
\InitSteps{} steps on the same data, batch and learning rate as the bound run,
and evaluated on the same \ToyN{} held-out pairs, giving \InitCells{} cells. An
initialisation probe recorded, before training, whether each cell's gates were
exactly zero.

The predeclaration fixed the stall criterion before any cell ran: a cell
stalls if its mean objective over the last ten steps is at least
\InitThreshold{}, and a mode stalls if at least two of its three seeds do. Four
reading rules were written in evaluation order, the first match to fire: no
cell stalls; every cell stalls; the default stalls and at least one alternative
escapes; anything else. Restored fidelity against the degraded input is
reported per cell but does not enter the rule, and no cell beyond the nine was
to be added whatever the outcome. The frozen sources were re-hashed after the
run and match; the repository head advanced during the run because of
unrelated commits, which the receipt records, and the code hash stamped in
every cell is identical. The default cell at seed zero is the bound
construction at the bound seed, and the figure code stops if it fails to
reproduce the bound objective and fidelity within the printed precision.

The scored means and the memory census are two populations, not one union. Let
\(\mathcal{U}\) be the candidate input pool. The displayed estimator is the
published RealSR crop128 protocol with a matched reference:
\[
  \mathcal{M}
  = \{i\in\mathcal{U}: i\text{ is in the published crop128 protocol and
      has a matched reference}\},
\]
with \(N_{\mathcal{M}}=|\mathcal{M}|=\SubsetN{}\). That protocol run recorded
\(\EvalSkipped{}\) memory skips and \(\EvalReferences{}\) matched references. A
separate earlier census applied a memory filter over \(\CensusTotal{}\) inputs;
write
\[
  \mathcal{C}
  = \{i\in\mathcal{U}: i\text{ was presented to that memory-filtered
      census}\}.
\]
The displayed estimator does not take \(\mathcal{M}\cup\mathcal{C}\) or
intersect \(\mathcal{M}\) with the memory filter. All subset means, paired
differences, intervals, and win counts are functions of \(\mathcal{M}\) only.
An input skipped for memory in the census is counted in \(\mathcal{C}\) and is
outside \(\mathcal{M}\); it contributes a census count, not a zero-valued
metric. Thus the reported comparison is the crop128 protocol and cannot be
read as an estimate for \(\mathcal{C}\) or for the released model on a full
benchmark.

\subsection{Checks the build enforces}

Several sentences in this paper are enforced as build-time checks rather than
asserted in prose, and the figure code stops rather than reporting a number that
would contradict them.

The one-line ladder entries that summarise the toy run must still restate the
values the training record holds. The record that withdrew the toy comparison
must agree with the training record on both parameter counts and with the
unpromoted longer run on its restored fidelity. Every mean printed in the subset
summary must be the mean of the per-crop rows it claims to summarise, and every
recorded interval must be centred on that same mean. The recorded paired
difference, its extremes, and the recorded win and loss counts, overall and
within each source group, must all reproduce from the per-crop rows. The per-crop
rows themselves are joined across arms on the crop name rather than on position,
so a reordering cannot silently pair the wrong crops.

Two recorded paths hold the record for the bicubic arm. Comparing their manifest bindings
establishes whether that is a second measurement or a second copy, and the
manuscript reports which; comparing their rows establishes that they still agree
either way.

\subsection{Bounding the model against itself}

Two pixel sets exist for the released model over the same \SubsetN{} crops,
produced by two inference passes and scored by the same protocol. The per-crop
differences between them are not used as a result. They are used once, to bound
how far re-running the recorded inference moved a crop. The record does not show
the sampling seed varied between the passes, so two runs cannot be read as an
estimate of sampling variance. The largest of the differences is drawn on the
figure that reports the gap, and it is reported as a largest observed shift
rather than as an interval.

\subsection{Tests}

The paired difference between the released model and bicubic upsampling is
tested three ways on the same \SubsetN{} crops: the exact two-sided sign
test~\cite{dixon1946sign}, which counts how many crops fall each way and assumes
nothing about their distribution; the Wilcoxon signed-rank
test~\cite{wilcoxon1945individual} on the same pairs; and the paired $t$
test~\cite{student1908probable} with its interval, which is the only one of the
three that assumes a distributional shape. A percentile
bootstrap~\cite{efron1979bootstrap} over the paired differences is reported
beside the interval each record already carries, so the two constructions can be
compared rather than one standing in for the other. Source group is a recorded
property of each crop and is reported as a stratum, not as a covariate: the
strata were not chosen after seeing the differences, and no adjustment is made.
The resampling unit is the crop, so these intervals describe this measured
subset and do not provide camera-level or population-level uncertainty when
crops from the same source are correlated.

For a crop \(i\in\mathcal{M}\), the fidelity contrast is
\(\Delta_i=\mathrm{PSNR}_{\mathrm{model},i}
-\mathrm{PSNR}_{\mathrm{bicubic},i}\) (and analogously for SSIM). These are
distortion contrasts, not latent perceptual-quality contrasts; a negative
\(\overline{\Delta}\) says only that the model is farther from the available
reference under that metric.

The cross-metric consistency panel computes this contrast separately for each
of RGB PSNR, luma PSNR, RGB SSIM and luma SSIM before taking the paired mean.
Each interval is the 2.5th and 97.5th percentiles of 1,000 resamples of the
corresponding paired differences. PSNR remains in decibels and SSIM remains
unitless, so the two metric families are displayed on separate native axes
rather than combined into a standardized score. This panel is descriptive: it
does not add a multiplicity-adjusted test or turn distortion differences into a
claim about perceptual quality.

\subsection{The census of an attempt, and the path it names}

An earlier attempt ran the same released model over the full set of
\CensusTotal{} inputs under a fixed memory budget, expressed as a maximum number
of output pixels and a maximum output side length. Inputs whose upsampled output
would exceed the budget were skipped before inference. The counts, the budget,
the size distribution of what was skipped and the source-group composition of
what was skipped are all recorded.

The census names a live path as its source, and that path \CensusPathState{}:
the later subset run overwrote it. A snapshot of the earlier state was retained
under a history directory, and the census is checked against that snapshot
rather than against the live path. The check is what makes the census usable;
without the retained snapshot the counts would be unverifiable, and the paper
reports the near miss rather than passing over it.

\subsection{The boundary of what is computed}

The perceptual computation is the single endpoint described above, and the
secondary backbone the protocol named is reported as not computed rather than
estimated. The \CensusRan{} outputs of the earlier attempt enter only as a
count: they have no reference images on disk, so they carry no value, and the
build stops if any of them ever acquires one, because their absence is the
reason those outputs cannot be scored. The two larger comparisons this line of
work planned stay at their recorded status --- \PrimaryTier{} for the
full-length reference comparison and \SotaTier{} for the comparison against the
strongest published method that would follow it --- and the subset is reported
as the subset, with completing it leaving both where they were. The scored
\SubsetN{} crops are the published crop128 protocol; the memory skip belongs to
the earlier census.
